\subsection{Tokens, componentes y reglas de composición}

Una interfaz de corrección reúne formularios, documentos, calificaciones y acciones
administrativas. Su coherencia requiere decisiones comunes sobre representación
y comportamiento. Como referencia de esta organización, el sistema de diseño de
la administración estadounidense, USWDS, describe los \textit{design tokens} como
elementos discretos de estilo que forman paletas consistentes de tipografía,
espaciado y color~\cite{codex_uswds_tokens}. Esta referencia se utiliza para
explicar el enfoque, sin implicar que SCDE incorpore la biblioteca USWDS.

En términos de arquitectura, conviene distinguir tres responsabilidades. Los
tokens asignan nombres reutilizables a valores visuales; las primitivas encapsulan
controles y estructuras comunes; las reglas de composición determinan cómo se
organizan esas piezas para realizar una tarea. La distinción permite revisar un
valor compartido sin buscarlo en cada pantalla, y estudiar el comportamiento de un
diálogo sin confundirlo con el formulario concreto que contiene. Esta separación
es una interpretación de diseño aplicada al proyecto, no un resultado experimental
sobre productividad.

En la base de SCDE inspeccionada, las variables CSS se concentran en
\texttt{src/styles/tokens.css} y las primitivas se exportan desde
\texttt{src/ui/index.jsx}. Las funcionalidades componen elementos como
\texttt{PageHeader}, \texttt{DataTable}, \texttt{Field} y \texttt{Dialog}, mientras
sus hojas locales resuelven la disposición del contenido. Por ejemplo, la hoja
administrativa emplea variables de espaciado y permite reorganizar los filtros
con una cuadrícula adaptable~\cite{codex_scde_design_evidence}. La utilidad de esta
organización reside en disponer de puntos concretos de revisión y modificación;
no demuestra por sí sola que todos los componentes satisfagan las necesidades
del alumnado o del profesorado.

\subsection{Accesibilidad como criterio verificable}

WCAG 2.2 organiza la accesibilidad mediante criterios comprobables, independientes
de una tecnología concreta. Para una herramienta de gestión docente resultan
pertinentes el contraste, el uso del teclado, el foco visible, las etiquetas y la
identificación de errores. La información de estado debe poder comprenderse sin
depender únicamente del color. Los niveles de conformidad exigen satisfacer los
criterios correspondientes en páginas y procesos completos; una comprobación
aislada no acredita la aplicación~\cite{codex_wcag22_ui}.

La densidad de una tabla también afecta a la precisión de las acciones. La
explicación oficial del criterio 2.5.8 sitúa el tamaño mínimo de los objetivos de
puntero en 24 por 24 píxeles CSS y contempla excepciones, entre ellas el
espaciado suficiente y los enlaces dentro de texto. Por ello, reducir el tamaño
de los botones para incluir más registros requiere examinar tanto su área
interactiva como la proximidad a otras acciones~\cite{codex_wcag_target_size}.
Esta observación orienta la revisión de tablas administrativas; aquí no se afirma
haber medido sus objetivos en un navegador.

Los diálogos muestran por qué la semántica y la interacción deben examinarse
juntas. El patrón modal de WAI-ARIA APG describe la entrada del foco en el
diálogo, su circulación interna con tabulación, el cierre mediante Escape y la
devolución del foco al elemento que lo abrió, con excepciones ligadas al flujo.
También especifica un nombre accesible y el estado modal. Para operaciones
difíciles de revertir aconseja considerar el foco inicial en la alternativa menos
destructiva~\cite{codex_apg_modal_dialog}. Estas indicaciones sirven como guía
de implementación y revisión, sin sustituir una evaluación del producto.

\subsection{Aplicación a SCDE y límites de la evidencia}

La guía v2 del proyecto prescribe una presentación clara y educativa, texto de
interfaz en Atkinson Hyperlegible~	te{atkinson_hyperlegible}, uso limitado de monoespaciada para códigos y
estados acompañados de significado textual. La guía, los \textit{tokens} claros
y las primitivas forman parte de la base compartida que ambas líneas de trabajo
heredan; durante el desarrollo coexistieron, no obstante, ramas con la versión
oscura anterior, de modo que las capturas y los resultados deben indicar siempre
la versión del código de la que proceden~\cite{codex_scde_design_evidence}.

La revisión de las primitivas desde la línea de administración permitió
detectar dos carencias concretas en una versión anterior de la base: la primitiva
\texttt{Dialog} enfocaba el diálogo, se cerraba con Escape y devolvía el foco al
elemento previo, pero no confinaba la tabulación; y \texttt{Field} asociaba la
etiqueta al control sin relacionar con él su ayuda o su mensaje de error. Ambas
se registraron como solicitudes de coordinación y se corrigieron en la base
compartida (commit \texttt{96fac49}): el diálogo mantiene ahora Tab y
Mayús+Tab dentro de sus elementos enfocables, y el campo vincula ayuda y error
mediante \texttt{aria-describedby} y marca \texttt{aria-invalid} cuando hay un
error. El episodio ilustra el valor de concentrar el comportamiento en
primitivas: una corrección en un único punto alcanza a todas las pantallas que
las usan. No sustituye, sin embargo, una auditoría de accesibilidad de la
aplicación completa.

El control automatizado de estilo detecta determinados patrones en los archivos,
como colores escritos directamente o modificaciones locales de clases comunes.
Las pruebas administrativas consultan controles por sus etiquetas y roles y
comprueban la confirmación de una desactivación. Ambas verificaciones aportan
evidencia acotada: ni el análisis textual ni la ejecución en un DOM simulado
reproducen toda la navegación con teclado, el contraste renderizado o el uso de
tecnologías de apoyo~\cite{codex_scde_design_evidence}. Una evaluación posterior
debe registrar versión, pantallas, tareas y dispositivos examinados. Así, el
sistema de diseño proporciona una base revisable, mientras la accesibilidad se
mantiene como una propiedad que requiere verificación explícita.
